\section{All algebraic orders of the signed radial sum}
\label{rad:section}

We first establish existence and a coefficient rule on the two level
parities, before assigning coefficients to Chern--Simons actions.
Beyond paths, the truncation walls at a branching vertex involve
signed sums of several neighbouring colours. Such walls can intersect
multiple endpoint poles. We neither replace them by an ordering of
colours nor discard nonzero aliases.

Let \(G=(V,E)\) be a fixed tree with \(N=|V|\ge2\), and put \(r=k+2\).
Use reflected one-based colours \(a_v=r-n_v\), \(1\le n_v<r\).
We write a radial test as \(F((n_v/r)_v)\). A smooth function of the
original fibre traces is of this form, with \(F\) extending smoothly
by even reflection at both endpoints of each coordinate. This
endpoint condition, more generally, suffices for our argument.
For \(d\ge1\) put \(W_d(x)=(1+\cos\pi x)/\sin^d(\pi x)\). For a fixed
parity vector \(\boldsymbol p\in\{0,1\}^V\) the signed kernel is
\begin{equation}
 \mathcal W_{\boldsymbol p,r}[F]=
 \sum_{\substack{1\le n_v<r\\n_v\equiv p_v\;(\mathrm{mod}\,2)}}
 F((n_v/r)_v)
 \prod_{v\in V}W_{\deg(v)}(n_v/r)
 \prod_{vv'\in E}\sin\!\left(\frac{\pi n_vn_{v'}}r\right).
 \label{rad:sum}
\end{equation}
The actual untwisted paired-cap amplitude is a finite signed linear
combination of these kernels times \((2/r)^{(N+1)/2}\).
The vacuum contribution uses only odd \(n_v\). A cap-product tag
\(q\) uses half the odd weight plus \((-1)^q\) times half the even
weight. Specifying an ordered pair of roots with that product
adds a factor \(1/2\). These are the exact cap identities in the
modular tree contraction \cite{GPPV}, not continuum replacements.
Reflection of an edge contributes \((-1)^{r+n_v+n_{v'}}\), fixed
on each parity choice. The untwisted surgery signature is zero:
the two opposite caps have cancelling Schur contributions, and the
unweighted bipartite backbone adjacency has spectrum symmetric
about zero.

\begin{theorem}[All-order radial expansion]\label{rad:thm}
For each fixed tree, each fixed smooth joint radial test, and each
fixed choice of physical cap projectors, the amplitude admits an
expansion on each parity \(r\equiv\rho\pmod2\). Its terms are
\(r^\alpha P_{\alpha,\rho,F}(\log r)\), where \(\alpha\) lies in a
fixed half-integral translate of \(\mathbb Z\), the power scale is
bounded above, and log degrees are bounded in terms of the tree.
For every \(J>0\), retaining all terms with \(\alpha\ge-J\) leaves
\(O(r^{-J})\); taking extra orders does not change these coefficients.
No other arithmetic subsequence is needed. The coefficients are
continuous linear functionals of finite order in \(F\). Finite
identities between the original signed marked sums hold termwise
in their expansions.
\end{theorem}

We prove this for \eqref{rad:sum}, then restore the prefactor and
signs. Multiplication by the ordinary asymptotic Taylor series of
\(S_{11}^{-1}\) gives the corresponding result for \(\tau_r=Z_r/S_{11}\).
Removing residual parity after the physical action is extracted
is a separate step.

\subsection{Character elimination: one affine wall per branch}
\label{rad:elimination}

Choose a bipartition \(A,B\) of \(G\). Expand \(F\) in its smooth
\(2\)-periodic Fourier series. By individual evenness we can use
the product of cosines at a mode \(\boldsymbol\sigma\in\mathbb Z^V\),
with the original joint Fourier coefficient
\(\widehat F(\boldsymbol\sigma)\). These coefficients decrease
rapidly. For any fixed \(\epsilon>0\), modes
\(\|\boldsymbol\sigma\|>r^\epsilon\) contribute arbitrarily small
powers of \(r^{-1}\), since the original finite sum in absolute
value is bounded by a fixed power of \(r\) for all unit-bounded tests.

At fixed \(A\)-colours and mode, the \(B\)-variables decouple.
At \(b\in B\) with \(d=\deg b\), its sines and denominator combine
into unnormalized characters, with continuous endpoint values:
\[
 \frac{\sin(\pi n_i m/r)}{\sin(\pi m/r)}
  =\sum_{j_i=0}^{n_i-1}
       \exp\!\left(\frac{i\pi(2j_i-n_i+1)m}r\right),
 \qquad i\in N(b).
\]
Put \(h_b(x)=(1+\cos\pi x)\cos(\pi\sigma_b x)\) and let \(h_{b,s}\)
be its coefficients in the basis \(e^{i\pi sx}\). Only
\(s=\pm\sigma_b,\pm\sigma_b\pm1\) can occur. Reflecting
\(m\mapsto 2r-m\) gives equal contributions. The possibly included
\(m=r\) term is zero, whereas the full-period even grid adds
\(m=0\). Exact finite orthogonality
\[
 \sum_{\substack{m\;(\mathrm{mod}\,2r)\\m\equiv p\;(\mathrm{mod}\,2)}}
 e^{i\pi\ell m/r}=r(-1)^{p\ell/r}{\bf1}_{r\mid\ell}
\]
with the right side interpreted as zero off \(r\mid\ell\), gives
the \(b\)-factor
\begin{equation}
 \begin{aligned}
 &\frac r2\sum_s h_{b,s}\sum_{\lambda\in\mathbb Z}
 (-1)^{p_b\lambda}{\bf1}_{L_b\in\mathbb Z}
 \sum_{U\subseteq N(b)}(-1)^{|U|}
 \binom{L_b-\sum_{i\in U}n_i+d-1}{d-1}_{+}\\
 &\quad-{\bf1}_{p_b=0}\frac{h_b(0)}{2}\prod_{i\sim b}n_i,
 \qquad L_b=\tfrac12\left(\lambda r-s+\sum_{i\sim b}(n_i-1)\right).
 \end{aligned}
 \label{rad:alias}
\end{equation}
The subscript \(+\) sets a binomial to zero if its upper argument
is below \(d-1\). Indeed the lattice condition requires
\(\sum_{i\sim b}j_i=L_b\), and excluding each event \(j_i\ge n_i\)
gives the displayed binomial. The final term subtracts the
artificial even-grid origin and must be kept for individual
Fourier modes, even if its aggregate vanishes. The contributing
aliases \(\lambda\) lie in a fixed finite interval if
\(|\sigma_v|\le r\), since each character frequency is bounded
by \(n_i-1\).

Expand the product of \eqref{rad:alias} as a signed sum. An
individual term selects one frequency, alias and subset, or
the origin correction, at each \(b\). Thus it keeps at most
\emph{one affine wall per \(b\)}: no common arrangement of all
competing walls is needed. Set \(r=2R+\rho\) and \(n_i=2j_i+p_i\)
for the retained labels. The integrality test for \(L_b\)
concerns only the fixed parities and the parity of \(s\).
For these allowed choices the binomial is a polynomial times
\begin{equation}
  \lambda_b R+\sum_{i\in N(b)}(1-2{\bf1}_{i\in U_b})j_i
              +c_b(\rho,\boldsymbol p,s)\ \ge0 ,
 \label{rad:wall}
\end{equation}
where the affine constant is integral. The other hard constraints
are \(n_i\ge1,\ r-n_i\ge1\). The remaining trigonometric factors
are smooth numerators times integer powers of those positive
endpoint distances, since \([x(1-x)]^dW_d(x)\) is smooth on
\([0,1]\). Polynomial weights and retained cosines multiply
this macro numerator. This reduction includes high-valence
walls as well as all signed aliases.

\subsection{A lattice decomposition with integral level step}
\label{rad:cone}

Fix the signs of the Fourier shifts and set
\(\upsilon_v=|\sigma_v|\). Homogenize \eqref{rad:wall}, the
endpoint bounds, and \(0\le\upsilon_v\le 2R+\rho\) in
\((R,(j_i)_{i\in A},(\upsilon_v)_{v\in V},H)\),
multiplying constants by \(H\). Add \(R\ge H\ge0\); the physical
slice has \(H=1\), and finitely many lowest levels are immaterial.
Impose the shift congruences on the integer lattice coset.
At \(R=0\) all variables vanish. The rational cone has compact
cross-section at \(R=1\). Pulling-triangulate this cross-section
with its original vertices and take the disjoint relative
interiors of all the resulting simplicial cones, including
lower-dimensional faces.

In one relative interior, a lattice point has a unique expression
\(u+\sum_i k_i e_i\) with independent integral ray vectors \(e_i\),
\(k_i\in\mathbb Z_{\ge0}\), and
\(u=\sum_i\theta_i e_i,\ 0<\theta_i\le1\). Scale rays if needed to
preserve the shift congruences. Only finitely many \(u\) occur.
At fixed \(H=1,\upsilon=|\boldsymbol\sigma|\), any \(k_i\) of a
ray carrying positive \(H\) or shift is \(O_G(Q)\), where
\(Q=1+\|\boldsymbol\sigma\|\), because these components are
nonnegative. Enumeration of those coefficients gives at most
polynomially many base points of size \(O_G(Q)\). Free rays have
\(H=\upsilon=0\).

Every free ray can be taken with \emph{unit \(R\)-component}.
Indeed, on the recession face \(H=\upsilon=0,\ R=1\), the
coefficient matrix in the \(j_i\) has entries \(\pm1\) supported
on a subforest of the original incidence graph, with at most
one row from each \(b\). A square minor cannot have two
contributing perfect matchings, whose symmetric difference
would form an alternating cycle. Consequently all square
determinants are \(0,\pm1\), and adjoining the endpoint unit
rows preserves this property. The right sides at \(R=1\) are
integral, so every recession vertex is integral by the determinant
formula for an independent subsystem of its active constraints.
As \(H=\upsilon=0\) is a face, these vertices are among the
original vertices of the full cross-section. The triangulation
adds no rays, and a free ray needs no congruence scaling since
its shifts vanish. This is where the forest and the single-wall
expansion avoid a new arithmetic period.

The lattice pieces at a fixed \(R\) are therefore simplexes
\begin{equation}
 k_i\in\mathbb Z_{\ge0},\qquad \sum_i k_i=R-R_0 .
 \label{rad:simplex}
\end{equation}
Pieces without a free ray occur only when \(r=O_G(Q)\), outside
the retained Fourier box for large \(r\). Each endpoint distance
is a positive affine form in the free \(k_i\), with base value
at least one and nonnegative slopes, by the recession box bounds.
On each simplex use a smooth partition of unity in barycentric
coordinates, choosing a pivot \(\ell\) with
\(k_\ell/(R-R_0)\ge1/(2m)\) on its support, where \(m\) is the
number of free rays. Eliminate \(k_\ell\) with
\eqref{rad:simplex}. The pivot support is a fixed normalized
distance from the omitted boundary and extends by zero; the
remaining lattice variables now range independently over
nonnegative integers. An endpoint distance seeing the pivot
with positive slope is \(\gtrsim r\) on the support for
\(r\gg_G Q\), so its inverse power gives a known power of \(r\)
and a smooth factor. Other endpoint forms retain their
nonnegative slopes.

Up to explicit powers of \(r\), the resulting sums are
\begin{equation}
 \sum_{k\in\mathbb Z_{\ge0}^d} f(tk,t)
       \prod_{j=1}^{D}\left(b_j+\sum_i a_{ji}k_i\right)^{-p_j}.
 \label{rad:orthant}
\end{equation}
Here \(t=r^{-1}\), \(b_j\ge1\), \(a_{ji}\ge0\), and
\(p_j\in\mathbb Z_{>0}\). The \(a_{ji},p_j\) belong to fixed finite families, and
\(b_j=O_G(Q)\). The profile has a smooth extension to a fixed
compact macro box for \(tQ\ll1\): the eliminated boundary is
avoided, the other orthant walls allow ordinary smooth extension,
and all offsets and retained frequencies are \(O_G(Q)\).
Every fixed derivative order thus costs at most a power of \(Q\).

\subsection{The positive affine orthant lemma}
\label{rad:orthantsec}

Smooth polytope Euler--Maclaurin formulas \cite{BV} and singular
lattice formulas for point-type weights \cite{BK} address related
settings. Here the needed estimate instead follows by resolving
a continuous Laplace kernel after the lattice sum is performed
exactly.

\begin{lemma}\label{rad:lemma}
The sum \eqref{rad:orthant} has an all-order expansion in integer
powers of \(t\) and polynomials in \(\log(1/t)\) of degree at most
\(D\). For each fixed output order both coefficients and remainder
bounds have at most polynomial cost in \(Q\). The power scale is
bounded below uniformly in the base and Fourier mode.
\end{lemma}

\begin{proof}
It suffices initially to take a \(t\)-independent profile \(f(x)\).
The general case follows by Taylor expansion in \(t\) with
polynomial-in-\(Q\) bounds in each specified derivative norm.
A remainder \(O(t^M Q^{C_M})\) costs only \(O(t^{-d})\) when summed
over the original compact support since each affine denominator
is at least one.

For \(g(x)=e^{\sum_i x_i}f(x)\), Fourier-invert its Schwartz
transform, Laplace-transform the denominator powers, and sum
the independent geometric series. The initial exchanges converge
absolutely for \(t>0\). Up to the Fourier convention, the sum is
the integral against \(\widehat g(\xi)\) of
\begin{equation}
 K(t,\xi)=\int_{y_j\ge0} e^{-\sum_j b_jy_j}
       \prod_j\frac{y_j^{p_j-1}}{\Gamma(p_j)}
       \prod_{i=1}^d\frac{1}{1-e^{-z_i}}\,dy .
 \label{rad:laplace}
\end{equation}
Here \(z_i=\sum_j a_{ji}y_j+t(1-i\xi_i)\). The real \(t\)-damping gives the convergent geometric sum and
prevents cancellation of its leading real scale. Since
\(\Re z_i\ge t\) and \(b_j\ge1\), we have
\(|K(t,\xi)|\lesssim (1-e^{-t})^{-d}\) uniformly in \(\xi\).
Insert a smooth cutoff \(\chi(t(1+|\xi|))\), equal to one near
zero and supported below \(c<1/2\). The change after Fourier
integration is \(O(t^M Q^{C_M})\) for any prescribed \(M\),
by the kernel bound and Schwartz decay.

Split the Laplace coordinates at one. In each region order all
small \(y_j\)'s and \(t\) as \(1\ge z'_1\ge\cdots\ge z'_L\),
ignoring null boundaries. Introduce ratios
\(z'_\ell=v_1\cdots v_\ell\), \(0\le v_\ell\le1\). If \(t\)
is in position \(u\), then \(t=\prod_{\ell\le u}v_\ell\).
Denominators seeing a large \(y_j\) with positive coefficient
have bounded smooth inverses with controlled derivatives.
For any other denominator let \(m(i)\) be the position of the
largest relevant variable, including \(t\). The quotient
\(z_i/z'_{m(i)}\) is smooth in the ratios with uniformly
positive real part, since one ratio is one with positive
coefficient. On the cutoff \(|\Im z_i|<c\), so
\(z/(1-e^{-z})\) is smooth, including at zero. It follows
that \(z'_{m(i)}/(1-e^{-z_i})\) extends smoothly to the
ratio cube, with derivative bounds polynomial in \(1+|\xi|\)
at every fixed order. Differentiating the Laplace weights
costs powers of \(Q\), and the large-coordinate integrals
are dominated by \(e^{-\sum_{\rm large}y_j}\) times polynomials.
The smooth cutoff allows extension by zero outside its region.

For the pushforward use \(dt/t\) and \(dy_j/y_j\) on the small
coordinates. The logarithmic change to the ratios has
determinant one. The monomial weight is
\(\prod_\ell v_\ell^{\beta_\ell}\), where
\begin{equation}
 \beta_\ell=\sum_{\substack{j\ {\rm small}\\{\rm pos}(j)\ge\ell}}p_j
    -\#\{i:m(i)\text{ exists and }m(i)\ge\ell\}.
 \label{rad:beta}
\end{equation}
The presence of \(t\) in every denominator gives \(m(i)\le u\).
Hence \(\beta_\ell>0\) for \(\ell>u\), and those ratios and
their required derivative bounds are integrable. Each sector's
Mellin transform in \(t\) is therefore
\begin{equation}
 \int_{[0,1]^u}\prod_{\ell=1}^{u}v_\ell^{w+\beta_\ell}
          A(v,\xi)\prod_{\ell=1}^{u}\frac{dv_\ell}{v_\ell}.
 \label{rad:mellin}
\end{equation}
Here \(A\) is smooth on the closed cube with the stated
seminorm bounds, and vanishes near the all-one corner by the
\(t\)-cutoff.

To control the remainder, apply a sufficiently long Taylor
subtraction at each \(v_\ell=0\) and expand the product of
Taylor and remainder operators. Polynomial terms supply
factors \((w+\beta_\ell+j)^{-1}\); the remaining integrals
are absolutely bounded on the prescribed strip. Poles thus
lie at integers and have order at most \(u\).
For vertical decay, split \(A\) into a finite smooth partition
on each piece of which some \(v_i\) is bounded strictly below
one. Such pieces cover its support. On the initial convergence
half-plane, repeated integration by parts in \(v_i\) with
\(v_i\partial_{v_i}\) has no boundary terms and divides each
time by \(w+\beta_i\). The differentiated amplitudes satisfy
the same Taylor/remainder strip estimates. Meromorphic
continuation gives arbitrary vertical decay on each strip
away from the poles with polynomial bounds in \(Q,1+|\xi|\)
at each fixed order.

Shift Mellin inversion to \(\Re w=-J-\tfrac12\). Residues
give integer powers of \(t\) and log-degree at most
\(u-1\le D\); the shifted integral gives \(O(t^{J+1/2})\)
with the polynomial parameter bounds. Only finitely many
poles are crossed, and \(\beta_\ell\ge-d\) gives the uniform
upper location bound on the poles. Integrating against
\(\widehat g\) preserves the expansion and its estimates.
\end{proof}

\subsection{Coefficient assembly and checks}
\label{rad:checks}

Apply Lemma~\ref{rad:lemma} to \eqref{rad:orthant} on
\(Q\le1+r^\epsilon\) with any fixed \(0<\epsilon<1\),
taking sufficiently many orders before restoring explicit
powers of \(r\). The number of base pieces and all fixed-order
bounds cost only polynomial powers of \(Q\), so the weighted
\(\widehat F(\boldsymbol\sigma)\) sum preserves the remainder.
The coefficient tails decay arbitrarily fast by the
polynomial bounds, and the omitted exact finite-sum tail
does so by the crude absolute bound on \eqref{rad:sum}.
Only finitely many smooth seminorms of \(F\) occur at a fixed
output order. This proves Theorem~\ref{rad:thm}.

A coefficient rule is now explicit: expand \eqref{rad:alias},
triangulate the affine cells, form \eqref{rad:laplace}, extract
the appropriate Mellin residues from \eqref{rad:mellin}, and
sum the Fourier coefficients absolutely. The exact partitions,
even-grid subtraction and signed inclusion--exclusion account
for the overlaps. Uniqueness of ordered power--log expansions
on each parity subsequence shows independence from the
triangulation, pivot partition and growing Fourier box.

\paragraph{Path check.}
For the unmarked four-vertex path, eliminating a leaf and a
degree-two vertex gives the exact expression
\begin{equation}
 Z_r(P_4)=(-1)^r\sqrt{\frac2r}
 \sum_{\substack{m,n<r\\ m,n\ {\rm positive\ odd}}}
 W_2(m/r)W_1(n/r)\min\{m,n,r-m,r-n\}.
 \label{rad:p4}
\end{equation}
The corresponding odd-grid \(W_1\) sine sum is \(r/2\);
the \(W_2\) two-sine sum is \(r/2\) times the displayed folded
minimum. These follow from \eqref{rad:alias}; for the latter
the zero and reflected aliases count the overlap of two
intervals of consecutive weights. The prefactor is
\((2/r)^{5/2}(r/2)^2=\sqrt{2/r}\). This checks the full alias
normalization, not a minimum-kernel hypothesis at higher
valence.

\paragraph{Intersecting branch walls.}
Let the retained set comprise \(a_0,\ldots,a_4\). The eliminated
vertices are two trivalent sites, with neighbours
\(\{a_0,a_1,a_2\}\) and \(\{a_0,a_3,a_4\}\), and an extra
leaf on \(a_0\). Both parts contain a trivalent site. In the
unaliased odd palette write \(n_i=2j_i+1\). If \(j_0\ge j_1+j_2\),
the first zero-weight count is
\((2j_1+1)(2j_2+1)\). In the adjacent regime
\[
 j_0=j_1+j_2-d,\qquad 0\le d\le 2\min(j_1,j_2),
\]
it decreases by \(d(d+1)\),
because each excluded signed tail has count \(1+\cdots+d\).
The second branch has the analogous wall \(j_0=j_3+j_4\).
They give simultaneous additive walls at the endpoint corner.
The shared \(a_0\) column is precisely the forest pattern
of Section~\ref{rad:cone}; outside the zero-weight window
the formula still includes all aliases. Neither part satisfies
a degree-two hypothesis, and no path estimate is needed at
the intersection.

The radial restriction is substantive. Angular Wilson
insertions can carry recoupling not represented by any
fixed smooth function of fibre traces; they need not reduce
to \eqref{rad:alias}. Backbone twists likewise insert
quadratic phases beyond this filter. This theorem also does
not assign different angular components with the same radial
projection their individual amplitudes.
